% GENERATED FILE. Edit canonical lessons, then run npm run generate:books.
% canonical-source-hash: fnv1a64:82f3b6d2307dd996
% canonical-lessons: BN-C40-words, BN-C40-lines, BN-C40-prothom-path

\chapter{Reading — Words, Lines, and a First Passage}
\label{ch:bn-prothom-path}
\hlchaptermodality{40}

\begin{chapteropening}
\textbf{By the end of this chapter:} \emph{I can read a Bengali word by its shape under the head-line, read the lines I would say to a stranger, and hold a whole passage without translating it word by word.}

Thirty-six words of Bengali with no new word in them, carried by three joiners of which only one is borrowed, and with the verb last in every line that has one.
\end{chapteropening}

\section[(reading)]{(reading) — six words under one line}

\label{lesson:BN-C40-words}



\begin{quote}
Every word below you have written. Not one is new.

What is new is that nobody says them first.
\end{quote}

\begin{tcolorbox}[breakable,skin=enhanced,colback=orange!6,colframe=orange!45!black,boxrule=0.5pt,arc=1mm,left=6pt,right=6pt,top=4pt,bottom=4pt,fonttitle=\bfseries,title={Read this}]
\begin{quote}
\bn{ভাত} \bn{চা} \bn{জল} \bn{দুধ} \bn{ভাই} \bn{বোন}
\end{quote}

Read down the list once, without stopping.

Again, and look at the top of each word rather than at the letters.
\end{tcolorbox}

\subsection*{What to know first}
Bengali hangs its letters from a line instead of standing them on one. That line runs unbroken across a whole word and stops at its edge.

So the line is a word boundary you can see before you have read anything. A Bengali word looks like one object first and a row of signs second, which is the opposite of how you learned to write it.

\textbf{\bn{ভাই}} and \textbf{\bn{বোন}} are the shortest pair here and the easiest to tell apart at a glance, for exactly that reason.

\subsection*{Your turn}
\begin{itemize}
\raggedright
  \item \emph{Say it:} the six words down the list, once, without stopping
\end{itemize}

\emph{Twice through:}

\begin{tcolorbox}[breakable,colback=teal!4,colframe=teal!35!black,title={Before you move on}]
How many of the six were new? (\textbf{None.})

What shows you where a Bengali word ends? (\textbf{The line stops.})
\end{tcolorbox}
\section[(reading)]{(reading) — six lines, and a bow nobody makes}

\label{lesson:BN-C40-lines}



\begin{quote}
The last lesson gave you words. These are things you would say.

You have said every one of them.
\end{quote}

\begin{tcolorbox}[breakable,skin=enhanced,colback=orange!6,colframe=orange!45!black,boxrule=0.5pt,arc=1mm,left=6pt,right=6pt,top=4pt,bottom=4pt,fonttitle=\bfseries,title={Read this}]
\begin{quote}
\bn{নমস্কার।} \bn{ধন্যবাদ।} \bn{হ্যাঁ।} \bn{না।} \bn{দয়া} \bn{করে।} \bn{স্বাগতম।}
\end{quote}

Read down once, without stopping.

Again — and notice the mark at the end of each line. It is not a full stop but a \textbf{\bn{দাঁড়ি}}, a single upright stroke.
\end{tcolorbox}

\subsection*{What to know first}
\textbf{\bn{নমস্কার}} is a bow written down: \emph{namas} plus \emph{kara}, the making of a bow. The commonest word in the language is a gesture that stopped being made.

\textbf{\bn{দয়া} \bn{করে}} is two words and means \emph{having done kindness}. Bengali asks for things by describing what the other person has already generously done, which is a habit worth noticing on the page where it sits complete.

Six lines, and four of them are a single word.

\subsection*{Your turn}
\begin{itemize}
\raggedright
  \item \emph{Say it:} the six lines, once, in order, without stopping
\end{itemize}

\emph{Twice through:}

\begin{tcolorbox}[breakable,colback=teal!4,colframe=teal!35!black,title={Before you move on}]
What ends a Bengali sentence? (\textbf{A \bn{দাঁড়ি}} — one upright stroke.)

What is \textbf{\bn{নমস্কার}}, literally? (\textbf{The making of a bow.})
\end{tcolorbox}
\section[(reading)]{(reading) — the first passage}

\label{lesson:BN-C40-prothom-path}



\begin{quote}
Six words, then six lines. Now thirteen of them together.

Thirty-six words, and not one of them is new.
\end{quote}

\begin{tcolorbox}[breakable,skin=enhanced,colback=orange!6,colframe=orange!45!black,boxrule=0.5pt,arc=1mm,left=6pt,right=6pt,top=4pt,bottom=4pt,fonttitle=\bfseries,title={Read this}]
\begin{quote}
\bn{নমস্কার।} \bn{আমার} \bn{নাম} \bn{রাম।} \bn{আমি} \bn{ভালো।} \bn{এটা} \bn{আমার} \bn{ভাই।} \bn{আমার} \bn{পরিবার} \bn{ভালো।} \bn{আমি} \bn{ভাত} \bn{আর} \bn{দুধ} \bn{খাই।} \bn{চা} \bn{বা} \bn{জল।} \bn{আমি} \bn{পড়ি।} \bn{আমার} \bn{বন্ধু} \bn{এখানে।} \bn{জামা} \bn{লাল}, \bn{কিন্তু} \bn{কাপড়} \bn{নীল।} \bn{প্রথম} \bn{আর} \bn{দ্বিতীয়।} \bn{হ্যাঁ}, \bn{ভালো।} \bn{ধন্যবাদ।}
\end{quote}

Read it once without stopping. Do not translate as you go — let the sentences arrive.

Now read it again and find the three joining words.
\end{tcolorbox}

\subsection*{What to know first}
\textbf{\bn{আর}} adds. \textbf{\bn{বা}} offers a choice. \textbf{\bn{কিন্তু}} turns the sentence.

Two of those are inherited and the third is not: \textbf{\bn{কিন্তু}} came through Persian, which is why it does not sound like its neighbours. You can hear the borrowing without being told about it.

Notice also that the verb is last in every line that has one. \textbf{\bn{আমি} \bn{ভাত} \bn{আর} \bn{দুধ} \bn{খাই}} — the eating comes after everything eaten, which is where Bengali always puts it.

This is the first whole passage the book has asked you to hold at once, and it works because everything in it was paid for.

\subsection*{Your turn}
\begin{itemize}
\raggedright
  \item \emph{Say it:} the passage aloud, once, without stopping
\end{itemize}

\emph{Twice through:}

\begin{tcolorbox}[breakable,colback=teal!4,colframe=teal!35!black,title={Before you move on}]
Where does a Bengali verb go? (\textbf{Last.})

Which of the three joiners is borrowed? (\textbf{\bn{কিন্তু}}, through Persian.)
\end{tcolorbox}
